% !TeX root = ../../Book5.tex
% 纲要标题：### 16.6 从代数空间到模问题的对象层次
% 纲要位置：几何学计划纲要.md，第 5675 行。
% 正文状态：未撰写；本文件目前是框架。

% 纲要内容开始
% 
% 本节为第33—34章建立共同语言；叠的实际定义和商叠模型在 V.34.5 展开，不提前要求导出叠或一般六函子理论。
% 
% * 集合值模函子与族的同构类
% * 概形可表和代数空间可表的不同结论
% * 保留自同构的群胚值对象及叠的动机
% * 通用族、刚化及粗模空间的区别
% * Hilbert、Picard 与商模问题的依赖清单
% 
% ---
% 
% 纲要内容结束

\section{从代数空间到模问题的对象层次}
\label{g5:sec:1606}

待撰写

\subsection{集合值模函子与族的同构类}
\label{g5:subsec:1606:01}
待撰写

\subsection{概形可表和代数空间可表的不同结论}
\label{g5:subsec:1606:02}
待撰写

\subsection{保留自同构的群胚值对象及叠的动机}
\label{g5:subsec:1606:03}
待撰写

\subsection{通用族、刚化及粗模空间的区别}
\label{g5:subsec:1606:04}
待撰写

\subsection{Hilbert、Picard 与商模问题的依赖清单}
\label{g5:subsec:1606:05}
待撰写
% 纲要条目排为编号小节标题，并进入目录与书签；说明段仍保留于纲要注释中。
